\subsection{Назначение программы и сценарии использования}
\label{subsec:program-purpose}

Программа предназначена для экспериментального исследования выполнения задач в режимах
  \texttt{seq}, \texttt{mn} и \texttt{pinned}, обоснованных в разделе~\ref{sec:threading-models}.
  Пользователь задаёт нагрузку и параметры исполнения, запускает обработку и получает результат
  вместе с длительностью выполнения. Программа представляет собой консольное приложение на Go,
  выполняемое на одной вычислительной системе. Сравнение платформ проводится посредством отдельных
  запусков на каждой из них.

Основной сценарий использования состоит в обработке одного подготовленного набора задач при
  выбранной конфигурации. При изменении режима сохраняются входные данные, полезные алгоритмы
  и состав нагрузки. В сериях измерений пользователь меняет один исследуемый фактор или заранее
  определённую группу параметров. Это позволяет сопоставлять результаты с гипотезами
  подраздела~\ref{subsec:research-hypotheses}.

Для расширения набора нагрузок предусматривается модульная организация. Подготовка сценария,
  реализация полезной работы, управление исполнением и вывод результатов имеют разные обязанности.
  Тип задачи определяет обработку данных, а исполнитель — порядок её запуска на горутинах.
  Добавление нового алгоритма должно сохранять общий механизм исполнения и измерения.

Единицей нагрузки является задача: обработка одного пакета записей или одного блока данных.
  Сценарий объединяет входные данные, список задач, их конфигурации и правила проверки результата.
  Один сценарий может содержать однотипные задачи либо смесь нескольких классов. Таким образом,
  число задач в сценарии не совпадает с числом горутин: одна горутина может последовательно
  обработать несколько элементов списка.

Требования к программе охватывают подготовку нагрузки, её параметризацию, три режима исполнения,
  проверку результата и сохранение сведений о запуске. Внешний источник входных данных не должен
  менять смысл конфигурации: параметры приводятся к единому внутреннему представлению и проверяются
  до начала обработки. Устройство компонентов и схемы их взаимодействия раскрываются
  в разделе~\ref{sec:results}.

% TODO: для текущей реализации конфигурация и сценарии задаются структурами Go в коде.
% YAML и фабрика задач относятся к следующей итерации; использовать те же структуры конфигурации.

\subsection{Типы задач, входные данные и результаты обработки}
\label{subsec:task-requirements}

Для выбранных в подразделе~\ref{subsec:workload-rationale} нагрузок предусматриваются два типа
  задач: \texttt{LogsTask} для обработки записей и \texttt{CompressTask} для сжатия блоков.
  Каждый тип представлен структурой с собственной конфигурацией и конструктором. Конструктор
  получает параметры и входные данные, проверяет их допустимость и создаёт экземпляр задачи.
  Выполнение полезной работы выделяется в метод \texttt{Run} с общим для всех типов контрактом.

Метод \texttt{Run} получает контекст управления запуском и возвращает результат либо ошибку.
  Контекст используется для передачи сигнала отмены. Настройки алгоритма хранятся в экземпляре
  задачи и во время обработки не меняются. Задача не выбирает режим многопоточности и не создаёт
  собственный пул горутин: эти действия принадлежат исполнителю. Такое разделение позволяет
  запускать один алгоритм последовательно и в двух параллельных режимах.

Один экземпляр задачи выполняется один раз. Повторный эксперимент получает новые экземпляры
  с теми же исходными параметрами и отдельным изменяемым состоянием. Входной набор допускает
  совместное чтение, а рабочие буферы принадлежат выполняющей обработку задаче. Общие изменяемые
  данные передаются явно как ресурс сценария и защищаются выбранным способом синхронизации.

Задача \texttt{LogsTask} получает пакет записей, каждая из которых содержит сервис, код ответа,
  длительность операции и размер сообщения. Она разбирает поля, применяет фильтр и рассчитывает
  статистику: число обработанных записей, суммы и гистограмму длительностей. Её параметры включают
  формат данных, число записей в пакете и условие фильтрации. Результат позволяет проверить
  полноту обработки и совпадение статистики с последовательным эталоном.

Длительности представляются целым числом в фиксированной единице времени. Счётчики и суммы
  накапливаются в целочисленных типах достаточной разрядности; допустимый объём сценария исключает
  переполнение. Границы интервалов гистограммы и правила включения граничных значений задаются
  конфигурацией. Число просмотренных записей учитывается до фильтрации; число принятых записей,
  суммы и гистограмма рассчитываются по записям, прошедшим фильтр. Средние значения вычисляются
  после объединения сумм и счётчиков.

Для накопления статистики предусматриваются две политики. При локальном подсчёте задачи возвращают
  частичные результаты, объединяемые в итоговый. При общей статистике обновления выполняются
  с защитой общим мьютексом. Обе политики должны приводить к одинаковым значениям. Политика
  накопления задаётся независимо от режима исполнения и сохраняется внутри сравниваемой пары.

Задача \texttt{CompressTask} получает блок байтов и конфигурацию сжатия. Результат содержит
  сжатые данные и сведения об обработанном объёме. Полнота преобразования проверяется восстановлением
  исходного блока и его сравнением с входом. Независимое сжатие блоков позволяет изменять число
  задач и их размер, сохраняя одну полезную операцию. Проверка не требует одинаковых сжатых байтов
  между разными программными платформами, если восстановленные данные совпадают с исходными.

Общий результат задачи содержит тип обработки, объём полезной работы и данные, необходимые для
  проверки. Время измеряется модулем эксперимента. Сам метод \texttt{Run} не записывает подробные
  логи по каждой операции, чтобы вывод не становился частью исследуемой нагрузки.

% TODO: сначала реализовать простые LogsTask и CompressTask; новые типы подключать к общему контракту.
% Поиск по ключу/диапазону и сетевой сценарий не включать в описание готовой реализации до их появления.

\subsection{Режимы исполнения и параметры запуска}
\label{subsec:execution-requirements}

Последовательный режим \texttt{seq} обрабатывает список задач без параллельной очереди.
  Режимы \texttt{mn} и \texttt{pinned} используют одинаковое число долгоживущих горутин,
  одинаковую очередь и одинаковые методы \texttt{Run}. В закреплённом режиме соответствие
  горутины потоку ОС устанавливается до сигнала готовности к обработке. После завершения
  закрепление снимается. Ограничения такого приближения модели 1:1 рассмотрены
  в подразделе~\ref{subsec:implementation-technologies}.

Общие параметры исполнения включают режим, значение \texttt{GOMAXPROCS}, число горутин,
  ёмкость очереди и объём входных данных. Эти величины задаются раздельно: число горутин определяет
  число исполнителей задач, а \texttt{GOMAXPROCS} — доступный среде Go параллелизм.
  Размер очереди ограничивает число ожидающих выдачи элементов и также сохраняется одинаковым
  в сравниваемых параллельных режимах.

Параметры нагрузки определяют размер пакета, количество записей, уровень сжатия и политику
  накопления результата. Разбиение должно покрывать весь входной набор, включая последний
  неполный блок. Размер блока задаётся вместе с единицей измерения: числом записей для разбора
  либо числом байтов для сжатия. Недопустимые значения выявляются до выделения основного объёма
  памяти и начала измерения.

Для сценариев с ожиданием и общим ресурсом предусматриваются отдельные параметры
  \texttt{wait\_duration} и \texttt{hold\_duration}. Первый задаёт ожидание без удержания
  общего мьютекса, второй — интервал его удержания в моделируемой критической секции.
  Ожидание таймера при удержании ресурса представляет специальный сценарий конкуренции;
  оно не используется для обозначения объёма вычислительной работы.

Длительность CPU-задачи определяется количеством полезных операций и свойствами данных.
  Завершение вычислений по истечении заданного времени привело бы к разному объёму работы
  на разных системах. Поэтому для вычислительной нагрузки конфигурируется объём, а длительность
  является результатом измерения. Для таймера задаётся целевой интервал ожидания; фактическое
  возобновление зависит от планирования и может произойти позже.

Параметры применяются на уровне сценария и отдельных задач. Это позволяет в одном списке
  использовать блоки разных размеров, сохраняя общий режим исполнения. При этом конфигурация
  должна однозначно определять каждую задачу, а все значения по умолчанию отражаются в сохранённых
  сведениях о запуске.

% TODO: варианты ожидания и удержания ресурса подключить после базового запуска двух типов задач.
% Конкретные интервалы и размеры для серий измерений закрепить в разделе 5 после пилота.

\subsection{Формирование нагрузки и воспроизводимость}
\label{subsec:reproducible-workloads}

Для исследования неоднородной нагрузки задачи объединяются в классы по объёму обработки.
  Лёгкая задача выполняет меньший объём, тяжёлая — больший при том же алгоритме. Например,
  два класса сжатия различаются размером блока, сохраняя уровень сжатия и правила подготовки
  содержимого. Это позволяет исследовать размер задачи отдельно от смены алгоритма.

Смешивание \texttt{LogsTask} и \texttt{CompressTask} задаёт другой вид сценария: одновременно
  меняется характер полезной работы. Его результаты рассматриваются отдельно от смеси размеров
  одного типа. Относительная стоимость классов устанавливается предварительными измерениями,
  поскольку название задачи не задаёт её фактическую длительность.

Состав смеси определяется точным количеством задач каждого класса. Доля тяжёлых задач указывается
  по числу элементов, а их доля в общем объёме данных учитывается дополнительно. Например,
  десять процентов крупных блоков могут содержать больше половины входных байтов.
  Если заданные доли не представимы при выбранном количестве задач без округления,
  сохраняются фактические количества и полученные доли.

Генерация использует два независимых начальных состояния: \texttt{data\_seed} для содержимого
  данных и \texttt{order\_seed} для порядка задач. Сначала создаётся набор с заданным составом,
  затем список перемешивается. Общий генератор случайных чисел внутри параллельной обработки
  не используется. Такой порядок исключает изменение входного набора из-за последовательности
  обращений горутин к генератору.

Каждая пара режимов получает одинаковые данные, конфигурации и список задач. Вместе с seed
  сохраняются сведения о генераторе, параметры состава и контрольные суммы данных и списка.
  Одного начального состояния недостаточно для воспроизведения нагрузки при изменении алгоритма
  генерации. Порядок выдачи фиксируется, но порядок начала и завершения обработки определяется
  планировщиками и не должен искусственно выравниваться между режимами.

Повторы с одним \texttt{order\_seed} показывают изменчивость исполнения одной нагрузки.
  Изменение этого параметра при фиксированных данных позволяет оценить влияние расположения
  длинных задач в очереди. Отдельный сценарий группирует тяжёлые задачи подряд, сохраняя их
  количество и входные блоки. Условия таких серий задаются до анализа результатов.

При сопоставлении смесей контролируется общий объём данных. Внутри пары режимов совпадают
  также число задач и их размеры. Если между разными составами смеси число задач меняется,
  это отражается в описании серии: изменение времени может включать влияние частоты обращений
  к очереди. Генерация и создание экземпляров задач выполняются до начала измеряемого участка.

\subsection{Нефункциональные требования}
\label{subsec:nonfunctional-requirements}

Корректность определяется совпадением полезного результата с эталоном, а не только успешным
  завершением процесса. Для обработки записей проверяются итоговые счётчики и гистограмма,
  для сжатия — восстановленные блоки. Идентификаторы задач позволяют установить, что каждая
  задача завершилась ровно один раз и никакой участок данных не был пропущен.

Программа должна исключать несогласованное изменение общих данных. Локальные результаты
  принадлежат отдельным задачам или горутинам, а доступ к общей статистике синхронизируется.
  Проверка гонок выполняется отдельно от измерений времени. Подготовка повторного запуска
  сбрасывает накопленные результаты и состояние ресурсов, чтобы предыдущая обработка
  не влияла на последующую.

Требование переносимости предусматривает сборку одной программной основы под x86-64 и arm64
  и нативное исполнение на выбранных системах. Результаты сопровождаются версией Go, ОС
  и параметрами сборки. Настройки среды не должны неявно менять заданный экспериментом
  доступный параллелизм. Ограничения на число системных потоков и память учитываются при
  выборе допустимой конфигурации.

Расход ресурсов должен быть контролируемым: число горутин и ёмкость очереди ограничены,
  общий объём данных известен до старта. Подготовленный список может хранить описания блоков
  и ссылки на общий неизменяемый набор, а не отдельную копию входа для каждой задачи.
  Буферы сжатия и сохранённые результаты также входят в расход памяти и должны учитываться
  при расчёте размера сценария.

Измеряемый интервал включает выдачу задач, полезную обработку и получение объединённого
  результата. Генерация данных, создание задач, подготовка горутин и проверка результата
  находятся вне него. Рабочие выделения памяти внутри \texttt{Run} учитываются как часть
  алгоритма. Подробный сбор диагностических данных отделяется от основной серии измерений.

При отмене программа прекращает выдачу новых задач, завершает участвующие горутины и освобождает
  удерживаемые ресурсы. Текущая обработка проверяет отмену в предусмотренных местах; немедленное
  прерывание произвольной операции не предполагается. Ошибка задачи не должна приводить
  к бесконечному ожиданию оставшихся элементов или к записи частичного результата как успешного.

Расширяемость обеспечивается общим контрактом задачи и единым представлением конфигурации.
  Новый тип обработки добавляется со своими параметрами и правилами проверки, сохраняя
  режимы исполнения и границы измеряемого участка. Подключение диагностики относится к управлению
  экспериментом и не требует изменения полезного алгоритма каждого типа задач.

% TODO: реализацию метрик, трассировок, дашбордов Grafana и анализа pprof оставить следующей итерации.
% Источник данных для Grafana и способ доставки определить отдельно. Диагностику запускать отдельно.
% При необходимости добавить схему контракта задачи; общую функциональную схему разместить в разделе 5.

\subsection{Сохраняемые результаты и обработка ошибок}
\label{subsec:output-requirements}

Результат запуска включает полезный выход программы и запись об эксперименте. Итоговый
  структурированный лог содержит идентификатор запуска, сценарий, режим, параметры нагрузки,
  число горутин, \texttt{GOMAXPROCS}, объём обработки, длительность и статус проверки.
  Для смеси дополнительно указываются фактический состав и оба seed. Сведения о платформе
  и программной среде связываются с записью запуска, чтобы сравнение можно было повторить.

Формат вывода должен допускать автоматическое чтение и построение сводных таблиц.
  Каждое числовое поле имеет определённую единицу: время, байты, записи или количество задач.
  Неизмеренный показатель не заменяется нулём. Если собираются память и число потоков ОС,
  фиксируются метод и момент сбора; эти значения отличаются от размера кучи Go и числа горутин.
  Исходные длительности сохраняются отдельно для каждого независимого повторения.

Запись итогового лога выполняется после завершения измеряемого участка. Подробный вывод
  по задачам не включается в основной режим. При отдельном измерении задержек событийные
  отметки собираются одинаковым способом в двух режимах; условия такого сбора отражаются
  в описании серии, поскольку дополнительные измерения могут влиять на обработку.

Проверка конфигурации различает неизвестный тип задачи, недопустимый размер блока, несовместимые
  параметры и невозможность подготовки указанного объёма данных. Ошибка выполнения содержит
  идентификатор задачи и этап возникновения. Ошибка записи результата также является ошибкой
  запуска: наличие вычисленного результата само по себе не подтверждает его сохранение.

Незавершённый запуск или несовпадение с эталоном фиксируются явно. Его длительность может
  сохраняться для диагностики, но не включается в сравнение успешных обработок.
  Такой порядок связывает требования к корректности с требованиями к измерениям: время
  имеет исследовательский смысл только при подтверждённом выполнении заданной полезной работы.
